%%% PACKAGES %%%
\usepackage[utf8]{inputenc}
\usepackage[T1]{fontenc}
\usepackage[french]{babel}
\usepackage{amsmath,amsfonts,amsthm,amssymb,amsbsy,amscd,mathrsfs}
\usepackage{lmodern}
%\usepackage{fourier}
\usepackage[scaled=0.875]{helvet} % ss
\usepackage{textcomp}
\usepackage{dsfont} % pour les symboles des ensembles
\usepackage[left=1.4cm,right=1.4cm,top=1.5cm,bottom=2cm]{geometry}
\usepackage{fancyhdr} 
\usepackage{multicol}
\usepackage{mdwlist}


%%% MISE EN PAGE %%%
\fancyhf{}
\setlength{\headheight}{26pt}
\setlength{\footskip}{18pt}
\renewcommand {\headrulewidth}{0pt}
\renewcommand {\footrulewidth}{0pt}
\usepackage{totpages}
\usepackage{makeidx}
\usepackage{graphicx} % inclure le fichier graphique de pondichery 2010
\usepackage{hyperxmp}
\usepackage[colorlinks=true,pdfstartview=FitV,linkcolor=blue,citecolor=blue,urlcolor=blue]{hyperref}
\pagestyle{fancy}
\usepackage{array}
\usepackage{tabularx}
\usepackage{multirow}
\usepackage{hhline}
\usepackage[table]{xcolor}
\usepackage{tikz}

\usepackage{mathtools} 		%pour l'alignement des termes d'une matrice
\usepackage{fltpoint}
\usepackage{xcolor,colortbl}
\usepackage{pstricks,pst-all,pst-plot,pstricks-add,pst-tree,pst-3dplot,pst-func}
\usepackage[french]{minitoc}
%\usepackage{ifthen}
\usepackage{calc} 			% Effectuer des calculs dans lateX
\usepackage[e]{esvect}		% Vecteurs \vv{u}
\usepackage{pifont}			% liste de dingbats
\usepackage{afterpage}
\usepackage[np]{numprint}	% Format nombres
\setlength{\parindent}{0mm}
\usepackage [alwaysadjust]{paralist}
\setdefaultenum {1.}{a)}{}{}
\hyphenpenalty=10000 
\tolerance=2000
\emergencystretch=1em
\widowpenalty=10000
\clubpenalty=10000
\raggedbottom
\renewcommand*{\tabularxcolumn}[1]{m{#1}} %centrage vertical des cellules d'un tableau tabularx

%%% MATHEMATIQUES
\newcommand{\euro}{\texteuro{}}
\newcommand{\R}{\mathds {R}}
\newcommand{\C}{\mathds {C}}
\newcommand{\N}{\mathds {N}}
\newcommand{\Z}{\mathds {Z}}
\newcommand{\Q}{\mathds {Q}}
\DeclareMathOperator{\e}{\mathrm{e}}
\renewcommand{\i}{\mathrm {i}}
\newcommand{\Ree}{\text{Re}}       % partie réelle
\newcommand{\Ima}{\text{Im}}       % partie imaginaire
\renewcommand{\mod}[1]{\left\| #1 \right\|}
\newcommand{\conj}[1]{\overline{#1}}
\newcommand{\abs}[1]{\left \lvert#1\right \rvert} % module

				% BASES VECTEURS $
\newcommand{\Oij}{\ensuremath{\left( {{\mathrm{O}};\vec{\imath},\vec{\jmath}} \right)}}
\newcommand{\Oijk}{\ensuremath{\left( {{\mathrm{O}};\vec i,\vec j ,\vec k} \right)}}
\newcommand{\OIJ}{$\left( O,I,J\right)$}             % Repère Oij
\newcommand{\Ouv}{$\left( O;\vv{u},\vv{v}\right)$}    % Repère Ouv
\newcommand{\VV}[3]{\vv{#1}\Coor{#2}{#3}}

				% MATRICES $
\newcommand{\Syst}[2]{\left\{\begin{array}{ccccc} #1\\ #2 \end{array}\right.} 
		% Systèmes 2X2 $\Syst{2x+1=0}{3x-1=1}$
\newcommand{\Coor}[2]{\begin{pmatrix} #1\\#2 \end{pmatrix}} 		% Coordonnées $\Coor{2}{3}$
\newcommand{\CoorE}[3]{\begin{pmatrix} #1\\#2\\#3 \end{pmatrix}} 	% Coordonnées $\CoorE{2}{3}{3}$	
%\begin{pmatrix}
%   a_1 & b_1 \\
%   a_2 & b_2 
%\end{pmatrix}

				% INTEGRALES $
\newcommand{\intx}[2]{\displaystyle{\int_{#1}^{#2}\,}}					% intégrales display
\newcommand{\prim}[3]{\bigl[#1 \bigr]_{#2}^{#3}}
\def\dx{{\,\text{d}x}}                   % Ecriture dx
\def\dt{{\,\text{d}t}}                   % Ecriture dt 
\def\dz{{\,\text{d}z}}                   % Ecriture dz
\newcommand{\dd}{\,\mathrm{d}}			% Ecriture d

				% LIMITES $
\newcommand{\lih}[2]{\displaystyle{\lim_{#1\to #2}\,}}	% limite général

				% SOMME $
\newcommand{\summ}[3]{\displaystyle{\sum_{#1 = #2}^{#3}}\,}				

				% FONCTIONS $
\newcommand{\fonction}[5]{\begin{array}{lrcl}
#1: & #2 & \longrightarrow & #3 \\
    & #4 & \longmapsto & #5 \end{array}}
	% {Nom fonction}{Df}{D_arrivée}{variable}{fonction}

				% INTERVALLE $
\newcommand{\oo}[2]{\bigl] \, #1 \, ; #2 \,\bigl[ }
\newcommand{\fo}[2]{\bigl[ \, #1 \, ; #2 \,\bigl[ }
\newcommand{\of}[2]{\bigl] \, #1 \, ; #2 \,\bigl] }
\newcommand{\ff}[2]{\bigl[ \, #1 \, ; #2 \,\bigl] }

\newcommand{\pa}[1]{\bigl( #1 \bigl)}	% parenthèses maths
\newcommand{\cro}[1]{\bigl[ #1 \bigl]}	% crochets maths
\newcommand{\alors}{\Longrightarrow}		% alors
\newcommand{\ssi}{\Longleftrightarrow}		% Si et seulement si
\newcommand{\nb}[1]{\nombre{#1}}

				% SUITES
\newcommand{\suite}[4]{\left\{\begin{array}{rcl}u_#1&=&#2\\u_{#3}&=&#4\end{array}\right.}
% $\suite{0}{124}{n+1}{u_n+3}$
% \begin{tabular}{rcl}$u_0$&& fixé,\\$u_{n+1}$&$=$&$u_n+r$\end{tabular}

				% Multilignes
%\begin{align*}
%y  = 1 + & \left(  \frac{1}{x} + \frac{1}{x^2} + \frac{1}{x^3} + \ldots \right. \\
%  &\left. \quad + \frac{1}{x^{n-1}} + \frac{1}{x^n} \right)
%\end{align*}


\DecimalMathComma
\newcommand{\fvalid}{\raisebox{1ex}{\rotatebox{180}{$\Rsh$}}} % flèche entrée pour prg casio


%%% CADRES
\newrgbcolor{Rfond}{.995 0.945 .985}  
\newrgbcolor{Rougef}{.8 0.1 .2}
\newrgbcolor{Rbord}{.6 0.1 .4}
\newrgbcolor{bleu}{0.1 0.05 .5}

\newcommand{\cadreR}[1]{%
\psframebox[framesep=.8em,linecolor=black,linewidth=0.25pt,framearc=0.1,fillstyle=solid,fillcolor=Rfond]{\begin{minipage}{\linewidth-1.6em}%
#1\end{minipage}}}
\newcommand{\cadre}[1]{%
\psframebox[framesep=.8em,linecolor=black,linewidth=0.25pt,framearc=0]{\begin{minipage}{\linewidth-1.6em}%
#1\end{minipage}}}

\newcommand{\encadrer}[1]{{\setlength{\fboxsep}{2mm}\fbox{#1}}}

%%% COMMANDES
\newcommand{\dem}{\medskip\par{\textsf{\textsc{{\small{\ding{106}\: démonstration}}}}}%
\medskip}

% Exemple
\newcounter{nexemple} 
\setcounter{nexemple}{0} 
\newcommand{\exemple}{
\refstepcounter{nexemple} 
\par{\textsf{\textsc{{\small{\ding{52}\: exemple \; \arabic{nexemple} \,:}}}}} \; %
\smallskip} 

% Propriété
\newcounter{nprop}[chapter]           % déclaration du numéro d'exo
\setcounter{nprop}{0}        % initialisation du numero
\newcommand{\prop}{%
\refstepcounter{nprop} 
\par{\textsf{\textbf{\textcolor{blue}{{\ding{171}\;\textsc{propriété \;  \arabic{nprop}\,:\quad}}}}}}
\medskip
}%
 
% Définition
\newcounter{ndefi}[chapter]           % déclaration du numéro d'exo
\setcounter{ndefi}{0}        % initialisation du numero
\newcommand{\defi}{%
\refstepcounter{ndefi} 
\par{\textsf{\textbf{\textcolor{blue}{{\ding{170}\;\textsc{définition \arabic{ndefi}\,:\quad}}}}}}
\medskip
}%

% Exercice_cours
\newcounter{nexer}[chapter]           % déclaration du numéro d'exo
\setcounter{nexer}{0}        % initialisation du numero
\newcommand{\exer}{%
\refstepcounter{nexer} 
\par{\textsf{\textbf{\textcolor{black}{{\ding{45}\;\textsc{exercice \arabic{nexer}\,:\;}}}}}}
\smallskip
}%

% Exercice_devoirs
\newcounter{nexo}[chapter]           % déclaration du numéro d'exo
\setcounter{nexo}{0}        % initialisation du numero
\newcommand{\exo}{%
\refstepcounter{nexo} 
\par{\textsf{\textbf{\textcolor{blue}{{\textsc{exercice  \; \arabic{nexo}\,:\quad}}}}}}
\bigskip
}%
   
% Remarques
\newcounter{nrem}
\setcounter{nrem}{0}
\newcommand{\rem}{
\refstepcounter{nrem} 
\par{\textsf{\textsc{{\small{$\blacktriangleright{}$\: remarque \; \arabic{nrem}\,:\quad}}}}} %
\smallskip}

% Méthode
\newcounter{nmeth} 
\setcounter{nmeth}{0} 
\newcommand{\meth}{
\refstepcounter{nmeth} 
\par{\textsf{\textbf{\textsc{\textcolor{black}{{ \small{méthode \; \arabic{nmeth} \,:}}}}}} }
\smallskip}



\usepackage{titlesec}

\titleformat{\part}[display]
{\sffamily\LARGE\bfseries\centering\scshape\Rbord}{\partname\ \thepart}{40pt}{}
\renewcommand{\chaptermark}[1]{\markboth{#1}{}}
\titleformat{\chapter}[display]
  {\sffamily\Large\bfseries\centering\Rougef}{\chaptertitlename\ \thechapter}{20pt}{\scshape\huge}

 

% en-tête
\lhead{\footnotesize{\textsf{Lycée VICTOR HUGO}}} %gauche
%\chead{\Rougef{\Large{\textsf{{\textbf{\textsc{\MakeLowercase{\leftmark}}}}}}}} %centre
%\rhead{\footnotesize{\textsf{T\up{le} ES 4}}\\ \footnotesize{\textsf{T\up{le} ES-L}}}%droit
%\rhead{\footnotesize{\textsf{BTS - SNIR1}} }%droit
%pied de page                                              
\cfoot{\scriptsize{\textsf{\thepage}}}
\lfoot{\scriptsize{\href{mailto:ricart.lvh@laposte.net}{ricart.lvh@laposte.net}}}
%\rfoot{\scriptsize{\href{https://padlet.com/ricartlvh/cours}{https://padlet.com/ricartlvh/cours}}}


%\lfoot{ \htmladdnormallink{\scriptsize{\textsf{\textsc{A. Yallouz (MATH@ES)}}}}{http://yallouz.arie.free.fr}} 

\fancypagestyle{plain}{%
\fancyhf{} % clear all header and footer fields
\rfoot{\scriptsize{\textsf{\thepage}}}
%\lfoot{ \htmladdnormallink{\scriptsize{\textsf{\textsc{A. Yallouz (MATH@ES)}}}}{http://yallouz.arie.free.fr}} 
\renewcommand{\headrulewidth}{0pt}
\renewcommand{\footrulewidth}{0pt}
}
 \makeatletter\@addtoreset{chapter}{part}\makeatother % renuméroter les chapitres pour chaque partie
\setcounter{secnumdepth}{3}
\makeatletter
\renewcommand\section{\@startsection {section}{1}{\z@}%
                                   {-3.5ex \@plus -1ex \@minus -.2ex}%
                                   {2.3ex \@plus.2ex}%
                                  {\reset@font\bfseries\red\sffamily\scshape\MakeLowercase}} 
                                 
\renewcommand\subsection{\@startsection {subsection}{2}{\z@}%
                                   {-3.5ex \@plus -1ex \@minus -.2ex}%
                                   {2.3ex \@plus.2ex}%
                                   {\reset@font\bfseries\blue\sffamily\scshape\MakeLowercase}}  
                                
\renewcommand\subsubsection{\@startsection {subsubsection}{3}{\z@}%
	{-3.5ex \@plus -1ex \@minus -.2ex}%
	{2.3ex \@plus .2ex}%
	{\reset@font\bfseries\sffamily\scshape\small\MakeLowercase}}
\makeatother
\setcounter{secnumdepth}{2}

\renewcommand\thesection{\Roman{section}}
\renewcommand\thesubsection{\arabic{subsection}}



%%%% debut macro %%%%
%\makeatletter
%\@addtoreset{nactivite}{section}
%\makeatother
%%%% fin macro %%%%
\newcommand{\refact}[1]{
\par{\textsf{\textbf{\textcolor{blue}{{\small{\textsc{activité}\refstepcounter{nactivite}\label{#1} \; \arabic{nactivite}}}}}}}\stepcounter{nactivite}} 

\renewcommand{\mtcSfont}{\small\sffamily \scshape}
\renewcommand{\mtcSSfont}{\footnotesize \sffamily}
\makeindex 

\setcounter{tocdepth}{2}

\newcommand{\clearemptydoublepage}{%
   \newpage{\pagestyle{empty}\cleardoublepage}}
   

\graphicspath{{./Images/}}

\newenvironment{liste}{%
   \begin{list}
      {$\bullet$}{%
         \setlength\itemindent{-0.5cm}
         \setlength{\itemsep}{0.1cm}}}
   {\end{list}}
   
\newcolumntype{C}[1]{>{\centering\arraybackslash}p{#1cm}}

\usepackage{listings}
\lstset{language=Python}
\usepackage{inconsolata}
\lstset{
    language=Python, %% Troque para PHP, C, Java, etc... bash é o padrão
    basicstyle=\ttfamily\small,
    numberstyle=\footnotesize,
    numbers=left,
    backgroundcolor=\color{gray!10},
    frame=single,
    tabsize=2,
    rulecolor=\color{black!30},
    %title=\lstname,
    escapeinside={\%*}{*)},
    breaklines=true,
    breakatwhitespace=true,
    framextopmargin=2pt,
    framexbottommargin=2pt,
    inputencoding=utf8,
    inputpath={{./scripts/}},
    extendedchars=true,
    literate=
  {á}{{\'a}}1 {é}{{\'e}}1 {í}{{\'i}}1 {ó}{{\'o}}1 {ú}{{\'u}}1
  {Á}{{\'A}}1 {É}{{\'E}}1 {Í}{{\'I}}1 {Ó}{{\'O}}1 {Ú}{{\'U}}1
  {à}{{\`a}}1 {è}{{\`e}}1 {ì}{{\`i}}1 {ò}{{\`o}}1 {ù}{{\`u}}1
  {À}{{\`A}}1 {È}{{\'E}}1 {Ì}{{\`I}}1 {Ò}{{\`O}}1 {Ù}{{\`U}}1
  {ä}{{\"a}}1 {ë}{{\"e}}1 {ï}{{\"i}}1 {ö}{{\"o}}1 {ü}{{\"u}}1
  {Ä}{{\"A}}1 {Ë}{{\"E}}1 {Ï}{{\"I}}1 {Ö}{{\"O}}1 {Ü}{{\"U}}1
  {â}{{\^a}}1 {ê}{{\^e}}1 {î}{{\^i}}1 {ô}{{\^o}}1 {û}{{\^u}}1
  {Â}{{\^A}}1 {Ê}{{\^E}}1 {Î}{{\^I}}1 {Ô}{{\^O}}1 {Û}{{\^U}}1
  {œ}{{\oe}}1 {Œ}{{\OE}}1 {æ}{{\ae}}1 {Æ}{{\AE}}1 {ß}{{\ss}}1
  {ű}{{\H{u}}}1 {Ű}{{\H{U}}}1 {ő}{{\H{o}}}1 {Ő}{{\H{O}}}1
  {ç}{{\c c}}1 {Ç}{{\c C}}1 {ø}{{\o}}1 {å}{{\r a}}1 {Å}{{\r A}}1
  {€}{{\EUR}}1 {£}{{\pounds}}1
}

\newcommand{\prog}[1]{
\begin{lstlisting}
#1
\end{lstlisting}
}
 % \lstinputlisting{./scripts/test.py} 

 %%%%%%%%%%%
 % EN test
 
 